% Part: first-order-logic
% Chapter: sequent-calculus
% Section: provability-quantifiers

\documentclass[../../../include/open-logic-section]{subfiles}

\begin{document}

\olfileid{fol}{seq}{qpr}
\olsection{\usetoken{S}{derivability}, పరిమాణీకరణలు}

\begin{explain}
  ఈ విభాగంలో స్థాపించే $\Proves$కు సంబంధించిన విషయాలు సీక్వెంట్ కలనం
  నియమాల నుంచి రావాలని కూడా సంపూర్ణతా సిద్ధాంతం కోరుతుంది.
\end{explain}

\begin{thm}
\ollabel{thm:strong-generalization} $c$ అనే స్థిరాంకం $\Gamma$లో గానీ
$!A(x)$లో గానీ కనిపించకుండా, $\Gamma \Proves !A(c)$ అయితే,
$\Gamma \Proves \lforall[x][!A(x)]$.
\end{thm}

\begin{proof}
ఏదో పరిమిత $\Gamma_0 \subseteq \Gamma$కు $\Gamma_0 \Sequent !A(c)$
యొక్క $\Log{LK}$-!!{derivation}ను $\pi_0$ అందాం. ఒక
$\RightR{\lforall}$ నిగమనాన్ని చేరిస్తే
$\Gamma_0 \Sequent \lforall[x][!A(x)]$ యొక్క !!a{derivation} వస్తుంది.
$c$ అనేది $\Gamma$లో గానీ $!A(x)$లో గానీ కనిపించదు కనుక ఐగెన్
చరరాశి షరతు తీరుతుంది.
\end{proof}

\begin{prop}
\ollabel{prop:provability-quantifiers}
\begin{tagenumerate}{prvEx,prvAll}
\tagitem{prvEx}{$!A(t) \Proves \lexists[x][!A(x)]$.}{}
\tagitem{prvAll}{$\lforall[x][!A(x)] \Proves !A(t)$.}{}
\end{tagenumerate}
\end{prop}

\begin{proof}
\begin{tagenumerate}{prvEx,prvAll}
\tagitem{prvEx}{$!A(t) \Sequent \lexists[x][!A(x)]$ సీక్వెంట్
  !!{derivable}:
  \begin{prooftree}
    \Axiom$!A(t) \fCenter !A(t)$
    \RightLabel{\RightR{\lexists}}
    \UnaryInf$!A(t) \fCenter \lexists[x][!A(x)]$
    \end{prooftree}}{}
\tagitem{prvAll}{$\lforall[x][!A(x)] \Sequent !A(t)$ సీక్వెంట్
  !!{derivable}:
  \begin{prooftree}
    \Axiom$!A(t) \fCenter !A(t)$
    \RightLabel{\LeftR{\lforall}}
    \UnaryInf$\lforall[x][!A(x)] \fCenter !A(t)$
    \end{prooftree}}{}
\end{tagenumerate}
\end{proof}

\end{document}
